\documentclass[11pt]{article}
\usepackage[a4paper,margin=1in]{geometry}
\usepackage{amsmath}
\usepackage{xcolor}

\definecolor{commentgray}{gray}{0.35}
\newenvironment{refcomment}%
  {\begin{quote}\color{commentgray}\itshape}%
  {\end{quote}}
\newcommand{\response}{\medskip\noindent\textbf{Response.}\ }
\newcommand{\point}[1]{\bigskip\noindent\textbf{#1}}

\setlength{\parskip}{0.4em}
\setlength{\parindent}{0pt}

\title{Response to Referee 2\\[2pt]
\large Manuscript JFM-2026-1781\\
\emph{Strain-coupled one-dimensional turbulence for rapid distortion}\\
Sharma \& Kerstein}
\date{}

\begin{document}
\maketitle

We thank the referee for the report and for the specific reading
recommendations, in particular the pointer to Sagaut \& Cambon (2018),
which has been adopted as the rapid-distortion reference spine. The
manuscript has been substantially rewritten: it is reorganised into three
declared parts around a single take-home message, Section~2 is compressed
and its two halves connected, redundant closure discussion is consolidated,
an independent DNS benchmark has been added, and the leading-edge
application flagged in the old conclusions is now carried out. Section,
figure and table numbers below refer to the revised manuscript.

\point{1. Length, readability, and the absence of a clear take-home
message.}
\begin{refcomment}
``Under its present form, this very long (42 pages) manuscript is very
difficult to read, with no clear take-home message(s) on its relevant and
original results. Often, similar discussions and models took place in
different sections (and plenty of sub- and subsubsections), with unuseful
redundancy.''
\end{refcomment}

\response
The paper is restructured into three explicitly announced parts (\S\,1,
closing paragraphs): the model and what a single line can know
(\S\S\,2--5); the measured limits and the relaxation mechanism (\S\,6); and
the consequence for leading-edge noise (\S\,7). The abstract and the
introduction now state the take-home message in one line---that the model
evolves the strained spectrum, carrying rapid distortion and nonlinear
relaxation together, and that the difference this makes to a downstream
(leading-edge-noise) prediction is a term too large to neglect. The
six contributions are listed once, as an itemised list in \S\,1. Redundant
discussion has been removed by moving the CMK comparison, the full
intervention catalogue, the implementation-verification detail and the
leading-edge-chain steps to electronic supplementary material (\S\S\,S1--S5)
and to Appendices~A--B, leaving each topic in a single place in the main
text.

\point{2. Section 2, and the disconnected subsections 2.1 and 2.2.}
\begin{refcomment}
``Section 2 illustrates what is the most difficult to understand \ldots\
with almost disconnected sections 2.1 and 2.2. \ldots\ Since the
application to flows with mean shear has been published to JFM in 2001 by
Kerstein et al., I wonder if a clear review of the JFM paper should be
better than the ms.\ subsection 2.1?''
\end{refcomment}

\response
Section~2 has been reorganised around a single organising principle---the
separation of the distortion physics by structure (production and a
spectrally closed rapid operator as continuous forcing; the eddy events as
the turbulence-driven relaxation)---which is stated at the outset and
carried through to the summary (\S\,2.6). Following the referee's
suggestion, the standard-ODT subsection (\S\,2.1) is now a compact review
keyed to Kerstein (1999) and Kerstein \textit{et~al.} (2001) rather than a
self-contained re-derivation. The link the referee found missing is made
explicit: \S\,2.2 derives the strain-coupled line equation and its
structural admissibility directly from the ODT kernel of \S\,2.1, and
\S\,2.3 connects the Reynolds-stress/RDT description to that line equation
through the rapid/slow pressure split, so 2.1 and 2.2 are now two stages of
one construction rather than two reviews.

\point{3. Sagaut \& Cambon (2018) as the reference set; RDT solution
machinery.}
\begin{refcomment}
``all details for solving RDT equations (2.14) by a tensorial kernel are
gathered in the monograph by Sagaut \& Cambon \ldots\ a new edition 2018
\ldots\ is of particular interest with its more complete chapter 8 on
turbulence with strain \ldots\ often more recent and more relevant that the
ones of the ms.''
\end{refcomment}

\response
Sagaut \& Cambon (2018) is now cited as the modern systematic reference for
rapid distortion in the introduction and is used as the reference for the
RDT solution machinery where the wavevector-ensemble integration is
introduced (\S\S\,1, 2.3). We have taken the referee's point that the exact
rapid pressure--strain and its spectral integration are standard results
and now present them as such, attributed to the classical sources and the
monograph, and separated from what is new here (the single-line reduction,
its bounds and the measured axisymmetry hypothesis in \S\,5).

\point{4. What are the objectives and achievements; the model description.}
\begin{refcomment}
``What are really the objectives and the achievements of the study remain
very difficult to size. \ldots\ Near the end of the paper, before Eq.\
(5.4) a useful possible description of the model is given: ODT evolves a
velocity vector $u'(y)$ on a single line aligned with the resolved
direction \ldots''
\end{refcomment}

\response
The objectives and achievements are now stated up front: the abstract gives
the aim and the outcome, and \S\,1 lists the six results with the section
where each is established and, in each case, the reference data it is
measured against. The concise description of the model the referee found
useful late in the paper has been moved forward: the line and its
observables are defined where the single-line question is set up (\S\,5.2,
equation~(5.3)), and the operational picture---ODT evolving $\boldsymbol
u'(y)$ on the resolved axis, with production and the rapid operator as
continuous forcing and the eddy events as relaxation---is now given in the
formulation and its summary (\S\S\,2.2, 2.6).

\point{5. The closure appears in many places.}
\begin{refcomment}
``The related problem of closure appears in different sections,
subsections, subsubsections \ldots\ p.\ 37, 5.3. Eqs.\ (5.7)-(5.10).
Perhaps the sole original result for the closure of $\Pi^{(r)}_{ij}$?''
\end{refcomment}

\response
The closure material is consolidated. The exact rapid term and the closure
problem are stated once in the formulation (\S\,2.4.1); the isotropic
moment closure actually used in the simulations is given once (\S\,2.4.2);
and the single-line treatment---which the referee identifies as the
original contribution---is collected in Section~5, where the collapsed
spectral form (now equations~(5.6)--(5.9)) and its kernels appear, with the
derivation in Appendix~A.3. Section~5.7 then shows explicitly that this
spectral form reduces to the moment closure used earlier, so the two
appearances of ``closure'' are the same object at two levels of resolution
rather than separate treatments. We agree with the referee's reading that
Section~5 is the original result, and the paper is now built around it.

\point{6. Validation, and the reliance on the 1985 Lee \& Reynolds DNS.}
\begin{refcomment}
``Validation of (several versions of) the model is not clear. Why the
author favours the old DNS by Lee and Reynolds (1985)? \ldots\ None of the
15 figures of the ms.\ gives a convincing validation of a new model,
capable of incorporating RDT and nonlinearity, vs.\ numerical or
experimental results.''
\end{refcomment}

\response
Validation has been strengthened and its versions delimited. Lee \&
Reynolds (1985) is retained because it resolves a Reynolds number a single
line can match and gives a graded moment-level target, but it is now one
condensed subsection (\S\,3.4) and its figures (4--6) carry the independent
DNS of Zusi \& Perot (2013, 2014) as an overlay, which agree where they
overlap. The convincing spectral validation the referee asks for is new:
figure~11 (with table~1) is a three-way comparison of the strain-induced
spectral anisotropy against our own $128^3$ strained-box DNS and its
exact-RDT companion at matched projection; figure~12 quantifies the
distance from linear theory versus strain rapidity; figure~16 extends this
to a five-way comparison across three rapidities; and figure~19 tests the
assembled leading-edge-noise chain against four published experiments. The
several ``versions'' are labelled explicitly (category.variant, table~4 and
supplementary \S\,S2) so that which model produced which figure is
unambiguous.

\point{7. Introduction references far from the model; shock-turbulence.}
\begin{refcomment}
``The introduction incorporates many references of studies of sound
radiation, that are very far from the model developped in the ms.\ Idem
with the shock-turbulence interaction \ldots''
\end{refcomment}

\response
The introduction has been rewritten and trimmed. The far-field
sound-radiation citations not connected to the distortion problem have been
removed, and the shock-turbulence material has been dropped; the retained
acoustics references are confined to the leading-edge-noise framework that
consumes the model output. The introduction now moves from the distortion
problem to the specific gap the model fills, without the surrounding survey.

\point{8. The conclusions open an application that is not carried out.}
\begin{refcomment}
``The last paragraph of the conclusions seems to open the way to an
application to complex flow: \ldots\ the distorted, anisotropic upwash
spectrum at the leading-edge plane \ldots\ I think that such an objective
merits to be supported, but this is not really done in the present ms.''
\end{refcomment}

\response
That application is now carried out in the body of the paper rather than
deferred. Section~7 supplies the computed distorted upwash spectrum to
Amiet's leading-edge-noise theory, both as a geometry-independent kernel
ratio (figure~17) and as an absolute far-field prediction at the reference
operating point (figure~18), quantifies the $5$--$20$\,dB difference from
the frozen-spectrum practice, and compares the assembled prediction with
four published experiments (figure~19), agreeing within $1$--$2$\,dB rms
where the facilities are cleanest. The conclusions now report this result
rather than proposing it.

\end{document}
